\subsection{Stage~K: Local Kubernetes Chain}
\label{sec:analysis-rq14}

Stage~K asks how end-to-end inference cost changes when ResNet-18 is
decomposed into synchronously chained coarse architectural stages under
in-cluster Kubernetes execution. The analytical contribution of this
subsection is not to repeat the per-condition latencies tabulated in
\cref{tab:rq14_results}, but to extract what the local Kubernetes
benchmark adds on top of the on-host findings of
\cref{sec:analysis-rq11,sec:analysis-rq12,sec:analysis-rq13} and what
it deliberately does not establish. Two analytical points carry the
subsection. First, the architectural ordering established on
\texttt{127.0.0.1} in Stage~L1 (more synchronously chained boundaries
imply more end-to-end cost in this ResNet-18 family~\cite{resnet})
remains visible once the workload is lifted into a Kubernetes pod
network, but the per-hop surcharge is not a constant. Second, the
chain overhead decomposes recognisably into a roughly linear platform
component and a tensor-payload-dependent component, in line with the
two-dimensional account synthesised in
\cref{sec:analysis-rq13}~\cite{dnnsplit,deeperthings}.

\paragraph{The architectural ordering survives the move into Kubernetes with measurable cost.}
Under the controlled local \texttt{kind} setup, mean end-to-end latency
increases monotonically from 81.066~ms for
\texttt{monolithic\_k8s\_1svc} to 83.387~ms, 87.728~ms, 90.572~ms, and
92.535~ms for \texttt{chain\_2svc} through \texttt{chain\_5svc}
(\cref{tab:rq14_results}). Relative to the Kubernetes monolithic
baseline, the overhead grows from +2.321~ms (+2.9\%) to +11.469~ms
(+14.1\%). The smallest non-monolithic overhead is therefore
\texttt{chain\_2svc} at +2.9\%, an empirical floor for this predefined
chain family rather than a carry-forward selection: as declared in
\cref{sec:methods-carryforward} and recorded in the frozen
\texttt{carry\_forward.json}, the carry-forward rule does not apply to
Stage~K because the Kubernetes condition set is predefined rather than
selected. The analytical reading is therefore continuity with Stage~L1
under a new platform: the same ResNet-18 chain family that showed a
monotonic split-cost gradient on \texttt{127.0.0.1} continues to do so
in-cluster, with a worst-case overhead bounded at +14.1\% on this
hardware.

\paragraph{The marginal cost of an additional boundary is shaped by the tensor moved at that boundary, not by a constant per-hop surcharge.}
The step from \texttt{chain\_4svc} to \texttt{chain\_5svc} adds
approximately 1.96~ms of mean latency, smaller than the preceding
2.84~ms step from \texttt{chain\_3svc} to \texttt{chain\_4svc}
(\cref{tab:rq14_results}). The recorded tensor-byte profile in the
frozen run explains the shrinkage: the runner's
\texttt{total\_activation\_bytes\_mean} rises from 1960~KiB to
2058~KiB across these two conditions, so the additional boundary that
distinguishes \texttt{chain\_5svc} from \texttt{chain\_4svc} transfers
the 98~KiB \texttt{layer4} output. This is the activation-transfer
dimension of \cref{sec:analysis-rq13} reappearing under
Kubernetes~\cite{deeperthings}: a later boundary that moves a small
tensor adds less than an earlier boundary that moves a larger one,
even when both are synchronous gRPC hops. The reading should be
stated narrowly: it is a property of the evaluated coarse-stage
ResNet-18 chain family under this local benchmark, not a universal
law about microservice depth.

\paragraph{The platform overhead is separable from compute and scales with hop count.}
\Cref{tab:rq14_overhead} reports the frozen runner's
\texttt{non\_compute\_overhead\_ms\_mean} alongside the per-condition
overhead and the recorded tensor totals. The non-compute / platform
column grows from 6.215~ms (\texttt{chain\_2svc}) through 9.435~ms,
12.088~ms, and 13.633~ms (\texttt{chain\_5svc}), while
\texttt{total\_compute\_ms\_mean} in
\texttt{merged\_results/condition\_summaries.csv} is approximately
flat across the chained conditions (77.17--78.90~ms versus 77.63~ms
for the monolithic baseline). Most of the added latency is therefore
not extra computation: it is per-request work attributable to the
synchronous gRPC hop and the surrounding Kubernetes plumbing
(kubelet, kube-proxy, the pod-network path, gRPC framing and
deserialisation). This pattern is consistent with the broader
characterisation of microservice RPC cost: layered gRPC / HTTP/2 /
TCP / IP stacks add measurable per-call overhead that is not absorbed
by collocation, and serialisation and kernel-network traversal are a
real share of microservice latency even between same-host
processes~\cite{rpc_overhead_hotos25,notnets_apsys24}. Two caveats
bound this comparison. The platform column is the runner's composite
\texttt{non\_compute\_overhead\_ms\_mean}; it is not decomposed into
kernel, kube-proxy, CNI, and gRPC sub-shares, and the Stage~K evidence
does not support such a decomposition. The cited papers measure the
RPC-tax share in different workloads (a key-value store and the
\textsc{DeathStarBench} hotel-reservation application), not chained
ResNet-18 inference, so they support the qualitative reading that
this column is non-trivial and grows with chain depth, not the
specific millisecond magnitudes observed here.

\paragraph{Scope of validity: single-node local Kubernetes.}
The deployment metadata records that for every condition, all service
pods and the benchmark client were scheduled on the single node
\texttt{thesis-rq14-control-plane} of the local \texttt{kind} cluster
(\texttt{merged\_results/deployment\_metadata.json}). The platform
overhead measured above therefore already includes the in-cluster
gRPC path and the kubelet-mediated pod network, but it does not
include a cross-node hop, a heterogeneous-node CNI traversal, or any
of the sources of variability that managed cloud Kubernetes
introduces. Cross-node sensitivity and managed-cloud transfer are
deferred to \cref{sec:analysis-rq15,sec:analysis-rq15b} respectively;
Stage~K itself should be read as a valid local in-cluster benchmark of
the predefined chain family, not as a statement about deployment
behaviour in production Kubernetes.

\paragraph{Validation.}
Three pieces of evidence support treating the within-run ordering and
the platform/compute decomposition as credible for this local
environment. First, functional parity is intact: both the local
segment-chain check and the live in-cluster gRPC chain check report
\texttt{max\_abs\_diff} \(=0.0\) at the predeclared tolerance of
\(1.0\times10^{-5}\) across all five conditions
(\texttt{parity\_validation.json}), so the chain latencies are not
confounded with numerical drift across the partitioned graph. Second,
round-level behaviour is stable: per-condition round CVs lie between
\(0.0074\) (\texttt{chain\_5svc}) and \(0.0161\)
(\texttt{chain\_3svc}) in \texttt{round\_consistency.json}
(\cref{tab:rq14_rounds}), in line with the reporting practice of
combining per-iteration and round-level statistics rather than mean
alone~\cite{benchmarking_sc15}. Third, infrastructure controls are
recorded and applied symmetrically across conditions: image
(\texttt{thesis-inference:latest}, pull policy \texttt{IfNotPresent}),
namespace (\texttt{rq14}), node, resource requests and limits
(1~CPU, 512~MiB per service pod; 1~CPU, 1~GiB client), threading,
CPU affinity to core~0, and \texttt{nice~-5} are identical across the
five conditions
(\texttt{merged\_results/deployment\_metadata.json},
\texttt{merged\_results/environment.json}), so the comparison is
internally fair on infrastructure. Three limits must be stated
explicitly. The image tag \texttt{thesis-inference:latest} is a
mutable Docker tag, so byte-exact reproduction of the local
\texttt{kind} run requires re-pulling and re-tagging the
corresponding ACR digest used in the AKS stages; the local result is
therefore reproducible only up to whatever image was loaded into the
kind cluster at run time (\cref{sec:methods-artifacts}). Governor
and turbo controls were requested but writes to \texttt{/sys} failed
because the containerised environment exposes those paths as
read-only; the detected governor was already \texttt{performance} and
turbo was already disabled, so the comparison is internally fair but
governor/turbo are observed rather than enforced. Finally, Stage~K is
a single host, single \texttt{kind} cluster, single-node placement,
fixed seed; cross-host or repeated clean-cluster replication is not
in scope for this stage.

\paragraph{Evaluation.}
The evaluative reading of Stage~K has three facets. First, against the
Kubernetes monolithic baseline, every chained condition is measurably
slower at this local setup, with relative overheads of +2.9\%,
+8.2\%, +11.7\%, and +14.1\% for \texttt{chain\_2svc} through
\texttt{chain\_5svc} (\cref{tab:rq14_results}) and effect sizes of
Cohen's \(d=0.614\), \(1.614\), \(2.392\), and \(2.874\) respectively
(\cref{tab:rq14_effects}). The Mann-Whitney \(p\)-values in
\cref{tab:rq14_effects} are near zero because the pooled
per-iteration sample size is large (\(N=1{,}000\) per condition); as
noted in the methodological note in \cref{sec:rq14}, they should not
be read as evidence of strong practical significance on their own,
and the round-level CVs and parity checks are the primary indicators
of result stability. Second, against a useful counterfactual, the
smallest non-monolithic overhead is \texttt{chain\_2svc} at
+2.321~ms (+2.9\%); within the predefined family it is therefore the
least expensive way to move from a single Kubernetes service to a
chained deployment while preserving exact functional parity in the
benchmarked workload. Third, against the boundaries of what Stage~K
can claim, the result is sufficient to support a local-in-cluster
feasibility claim --- the chained ResNet-18 family runs end-to-end,
with bounded overhead and without functional divergence from the
monolithic baseline --- but it is not sufficient to claim production
readiness, scalability, multi-tenant robustness, or portability of
the absolute millisecond values. Cross-node sensitivity,
managed-cloud transfer, and service-mesh hardening are addressed in
\cref{sec:analysis-rq15,sec:analysis-rq15b} and in the security
analysis subsections respectively; Stage~K's contribution stops at the
local-feasibility boundary.

\paragraph{Stage~K outcome.}
Under the controlled local Kubernetes setup, increasing the number of
synchronously chained coarse-stage ResNet-18 services raises
end-to-end latency from \texttt{monolithic\_k8s\_1svc} through
\texttt{chain\_5svc}, with mean latency increasing from 81.066~ms to
83.387~ms, 87.728~ms, 90.572~ms, and 92.535~ms. The smallest
non-monolithic overhead is \texttt{chain\_2svc} at +2.321~ms
(+2.9\%), and the worst-case overhead is \texttt{chain\_5svc} at
+11.469~ms (+14.1\%). The marginal cost of an added boundary is not a
constant per-hop surcharge: the final increment from
\texttt{chain\_4svc} to \texttt{chain\_5svc} is approximately
1.96~ms, smaller than the preceding 2.84~ms step, because the
additional boundary moves only the 98~KiB \texttt{layer4} output.
Most of the added latency is platform-side rather than compute-side:
total compute is approximately flat across the chained conditions
while the recorded non-compute / platform overhead grows from
6.215~ms to 13.633~ms, consistent with the measurable per-hop cost
of layered gRPC microservice
communication~\cite{rpc_overhead_hotos25,notnets_apsys24}.
Carry-forward is not applicable in this stage because the Kubernetes
condition set was predefined rather than selected. These claims are
scoped to a single-node local \texttt{kind} cluster with all pods
colocated on \texttt{thesis-rq14-control-plane}; whether the
ordering and the bounded marginal cost transfer to managed cloud
Kubernetes and to multi-node placement is the empirical question
that \cref{sec:analysis-rq14b,sec:analysis-rq15,sec:analysis-rq15b}
address.
